\documentclass[11pt]{article}

\usepackage[a4paper,margin=1in]{geometry}
\usepackage[T1]{fontenc}
\usepackage[utf8]{inputenc}
\usepackage{lmodern}
\usepackage{amsmath,amssymb,amsfonts}
\usepackage{bm}
\usepackage{mathtools}
\usepackage{hyperref}
\usepackage{enumitem}
\usepackage{xcolor}
\usepackage{microtype}

\hypersetup{
  colorlinks=true,
  linkcolor=blue!50!black,
  citecolor=blue!50!black,
  urlcolor=blue!50!black
}

\title{Static/Dynamic Regions Estimation Across Video Frames}
\author{}
\date{}

\begin{document}
\maketitle

\begin{abstract}
We describe a unified per-camera front-end for estimating static and dynamic image regions across video frames. The framework combines dense point tracking, motion prediction, memory-based temporal association, camera-compensated motion analysis, and bidirectional region fusion. The resulting output is a set of persistent pixel-space tracklets together with a temporally coherent dynamic mask for each frame. The method is a synthesis of ideas drawn from prior work rather than a direct reproduction of a single existing approach: dense point tracking is inspired by TAPIR~\cite{tapir2023}, temporal memory by SAM~2~\cite{sam2_2024}, position-aware propagation by SAMITE~\cite{samite2025}, motion-guided prediction by SORT~\cite{sort2016} and SAMURAI~\cite{samurai2024}, the static-versus-dynamic criterion by MotionGS~\cite{motiongs2024}, and temporal region fusion by DEVA~\cite{deva2023}.
\end{abstract}

\section{Problem formulation}
Let $I_t$ denote the image at frame $t$ for one camera, with $t \in \{1,\dots,T\}$. The goal is to estimate, for each frame, a soft dynamic mask $M_t \in [0,1]^{H\times W}$, whose complement $1-M_t$ corresponds to the static region. At the same time, we seek to maintain persistent pixel-space trajectories so that dynamic evidence is accumulated over time rather than inferred independently at each frame.

We represent the temporal state of the image with a set of persistent seed tracks
\begin{equation}
\mathcal{T}=\{\tau_i\},
\qquad
\tau_i = \{x_{i,t},\, \hat{x}_{i,t|t-1},\, d_{i,t},\, v_{i,t},\, m_{i,t}\}_{t=1}^{T},
\end{equation}
where $x_{i,t} \in \mathbb{R}^2$ is the observed image position of seed $i$, $\hat{x}_{i,t|t-1}$ is its predicted location prior to association, $d_{i,t}$ is a descriptor or memory embedding, $v_{i,t} \in [0,1]$ is a visibility/confidence score, and $m_{i,t} \in [0,1]$ is the probability that the seed belongs to a dynamic region.

The main design principle is the following:
\begin{quote}
A pixel should be classified as \emph{dynamic} only if its observed motion cannot be explained by camera motion and the static scene hypothesis.
\end{quote}
This criterion is directly inspired by the motion decomposition viewpoint of MotionGS~\cite{motiongs2024}.

\section{Dense seed initialization and point tracking}
We first initialize a dense but manageable set of seeds in each frame using a hybrid strategy: a regular image grid ensures broad coverage, additional seeds are placed on high-gradient or corner-rich locations, and new seeds are spawned in newly visible regions when track density drops below a desired threshold.

Each seed is then tracked through time using a TAPIR-style point tracking backbone~\cite{tapir2023}. TAPIR is particularly relevant here because it treats the primitive as a persistent queried point and performs tracking through a combination of point matching and temporal refinement. In our formulation, the tracking backbone provides a measurement proposal
\begin{equation}
 z^{\mathrm{tap}}_{i,t+1}=\mathrm{TrackPoint}(x_{i,t}, I_t, I_{t+1}, d_{i,t}),
\end{equation}
which is later combined with a motion prior and temporal memory.

The reason for using a point tracker as the basic temporal primitive is that dynamic/static classification is far more stable when attached to persistent tracks than when estimated independently at the pixel level in each frame.

\section{Prediction before association}
To avoid full-frame matching at every step and to stabilize temporal linking, each seed maintains a small image-plane motion state,
\begin{equation}
 s_{i,t} = [u_{i,t}, v_{i,t}, \dot{u}_{i,t}, \dot{v}_{i,t}]^{\top}.
\end{equation}
We propagate this state with a linear motion model,
\begin{equation}
 \hat{s}_{i,t+1|t} = A s_{i,t|t},
 \qquad
 \hat{x}_{i,t+1|t} = H \hat{s}_{i,t+1|t},
\end{equation}
where $A$ is the transition matrix and $H$ extracts the predicted position component.

This predict-then-associate structure is inspired by classical tracking-by-detection, especially SORT~\cite{sort2016}, which uses a Kalman filter before assignment. In the context of SAM-based video tracking, SAMURAI~\cite{samurai2024} reinforces the value of this design by incorporating motion-aware memory selection and motion-guided prediction to improve temporal propagation. In our framework, the motion model is not meant to be a full tracker by itself; instead, it provides a local positional prior that constrains the subsequent association step.

We therefore define a gated search region
\begin{equation}
 \Omega_{i,t+1} = \{x \in \mathbb{R}^2 : \|x-\hat{x}_{i,t+1|t}\|_2 < r_i\},
\end{equation}
where $r_i$ is an adaptive search radius. Restricting association to $\Omega_{i,t+1}$ reduces drift, suppresses distractors, and allows tracks to survive short periods of ambiguity or occlusion.

\section{Memory-based temporal association}
Prediction alone is insufficient in the presence of repeated patterns, occlusions, appearance changes, and local ambiguities. For this reason, each seed is also associated with a memory bank
\begin{equation}
 \mathcal{M}_i = \{(d_{i,k}, x_{i,k}, q_{i,k})\}_{k\in\mathcal{H}_i},
\end{equation}
where $d_{i,k}$ is a stored descriptor, $x_{i,k}$ is the corresponding image location, $q_{i,k}$ is a reliability score, and $\mathcal{H}_i$ indexes selected historical observations.

The use of memory is motivated by SAM~2~\cite{sam2_2024}, whose video mechanism relies on streaming memory to condition current-frame prediction on past frames. However, SAM~2 does not explicitly impose pointwise spatial correspondence for dense pixel tracking. To introduce spatial guidance, we also borrow the position-aware idea of SAMITE~\cite{samite2025}, which augments SAM~2 with a calibrated memory bank and a positional prompt generator.

Inside the gated search region, candidate associations are scored as
\begin{equation}
 S_i(x)=
 \lambda_a\,\mathrm{sim}(d_{i,t}, d_{t+1}(x))
 +\lambda_m\,\mathrm{MemScore}_i(x)
 +\lambda_p\,\exp\!\left(-\frac{\|x-\hat{x}_{i,t+1|t}\|_2^2}{2\sigma_i^2}\right),
\end{equation}
where the first term measures descriptor similarity, the second aggregates consistency with the seed memory bank, and the third is a positional prior centered at the predicted location. The resulting associated observation is denoted by
\begin{equation}
 x^{\mathrm{obs}}_{i,t+1} = \arg\max_{x \in \Omega_{i,t+1}} S_i(x).
\end{equation}

After association, the seed state can be updated in Kalman-filter form,
\begin{equation}
 s_{i,t+1|t+1} = \hat{s}_{i,t+1|t} + K_{t+1}\big(z_{i,t+1}-H\hat{s}_{i,t+1|t}\big),
\end{equation}
where $z_{i,t+1}$ is the selected measurement and $K_{t+1}$ is the Kalman gain. Even when the eventual implementation departs from a strict linear-Gaussian filter, this predict-correct abstraction remains a useful conceptual layer for the front-end.

\section{Static-motion explanation and dynamic residual}
The central step of the framework is to estimate how each tracked seed would move under the \emph{static-scene hypothesis}. Let
\begin{equation}
 x^{\mathrm{stat}}_{i,t+1}
\end{equation}
be the predicted next position of seed $i$ if the scene were static and only the camera moved. This can be estimated in two ways:
\begin{enumerate}[leftmargin=1.5em]
 \item if camera poses and depth are available, by warping the point under ego-motion;
 \item otherwise, by fitting a robust epipolar geometry model from tracks that are currently most consistent with static motion.
\end{enumerate}

The dynamic residual is then defined as
\begin{equation}
 r_{i,t} = \|x^{\mathrm{obs}}_{i,t+1} - x^{\mathrm{stat}}_{i,t+1}\|_2.
\end{equation}
This is the core static-vs-dynamic discriminator. The idea comes directly from MotionGS~\cite{motiongs2024}, which explicitly separates observed motion into camera-induced flow and motion flow. In our image-space setting, a small residual indicates that the observed motion is well explained by the static scene, whereas a large residual suggests genuine dynamic behavior.

\section{Dynamic probability of each track}
We convert the residual into a soft dynamic probability by combining it with association quality and track confidence,
\begin{equation}
 m_{i,t+1}
 =
 \sigma\!\Big(
 \alpha r_{i,t}
 +\beta(1-v_{i,t+1})
 +\gamma(1-\mathrm{AssocConf}_{i,t+1})
 +\delta\,\mathrm{AppearanceChange}_{i,t}
 \Big),
\end{equation}
where $\sigma(\cdot)$ is the logistic sigmoid, and $\alpha,\beta,\gamma,\delta$ are weighting coefficients.

To reduce flicker, we smooth the dynamic probability over time,
\begin{equation}
 \bar{m}_{i,t+1} = \rho\,\bar{m}_{i,t} + (1-\rho)\,m_{i,t+1},
\end{equation}
and apply hysteresis:
\begin{equation}
 \tau_i \text{ is dynamic if } \bar{m}_{i,t} > \tau_{\mathrm{on}} \text{ for } K \text{ consecutive frames},
\end{equation}
while a dynamic track returns to static only when $\bar{m}_{i,t} < \tau_{\mathrm{off}}$ with $\tau_{\mathrm{off}} < \tau_{\mathrm{on}}$.

This step is a synthesis rather than a direct copy of a single paper. The residual-motion criterion comes from MotionGS~\cite{motiongs2024}; the confidence-aware formulation is naturally supported by point-tracking and memory-based association; and the temporal persistence reflects the same motivation as temporally coherent video segmentation approaches.

\section{From dynamic tracklets to dynamic regions}
The previous stages classify \emph{tracks}, not dense image regions. To obtain a per-frame dynamic mask, we first convert dynamic tracklets into a seed mask $M_t^{\mathrm{seed}}$, for example by splatting track confidence into local neighborhoods or by assigning the dynamic label to the superpixels containing highly dynamic tracks.

This initial mask is then refined by temporal propagation and fusion. Here we adopt the main design principle of DEVA~\cite{deva2023}: rather than trusting each frame independently, we fuse dynamic evidence from the current frame with information propagated from temporally adjacent frames. Specifically, we define
\begin{equation}
 M_t^{\star} = \mathrm{Fuse}\!\left(
 M_t^{\mathrm{seed}},
 \overrightarrow{M}_{t-k\rightarrow t},
 \overleftarrow{M}_{t+k\rightarrow t}
 \right),
\end{equation}
where $\overrightarrow{M}_{t-k\rightarrow t}$ and $\overleftarrow{M}_{t+k\rightarrow t}$ denote masks propagated forward and backward from neighboring frames. The fusion function may be implemented with confidence-weighted averaging, recurrent refinement, or graph-based region optimization.

The role of DEVA in our framework is therefore not semantic segmentation itself, but the \emph{bidirectional temporal fusion principle} that turns sparse and noisy dynamic evidence into temporally coherent region estimates~\cite{deva2023}.

\section{Unified algorithmic view}
Putting the pieces together, the per-camera front-end proceeds as follows for each transition $t \rightarrow t+1$:
\begin{enumerate}[leftmargin=1.5em]
 \item predict the next image location of every seed using the motion state update (SORT/SAMURAI-inspired)~\cite{sort2016,samurai2024};
 \item associate the seed locally with TAPIR-style point tracking constrained by a SAM~2/SAMITE-inspired memory-and-position prior~\cite{tapir2023,sam2_2024,samite2025};
 \item estimate the static-scene motion explanation and compute the camera-compensated residual (MotionGS-inspired)~\cite{motiongs2024};
 \item update the dynamic probability of each track and decide whether it belongs to a static or dynamic region;
 \item convert dynamic tracks into region seeds and fuse them over time using bidirectional propagation (DEVA-inspired)~\cite{deva2023};
 \item update memory banks, visibility scores, and track states for the next frame.
\end{enumerate}

In compact form, the processing loop is
\begin{equation}
 \text{Predict} \rightarrow \text{Associate} \rightarrow \text{Explain by static motion} \rightarrow \text{Score dynamicity} \rightarrow \text{Fuse masks} \rightarrow \text{Update memory}.
\end{equation}

\section{Discussion}
The proposed framework is intentionally modular. Different point trackers (e.g., TAPIR or CoTracker-style alternatives), different motion models, and different mask-fusion implementations can be substituted without changing the core logic. The main conceptual contribution of the framework is the integration of five ideas into a single front-end:
\begin{enumerate}[leftmargin=1.5em]
 \item persistent pixel-space trajectories as the temporal primitive;
 \item prediction before association to avoid full-frame search;
 \item memory-based matching with explicit positional priors;
 \item dynamicity defined as motion unexplained by the static-scene hypothesis;
 \item temporally fused region estimation rather than framewise masking.
\end{enumerate}

This front-end is particularly suitable as the first stage of a later Gaussian-based dynamic scene model, because it provides exactly the two pieces of information needed upstream: persistent image-space tracks and temporally coherent dynamic/static region masks.

\begin{thebibliography}{99}

\bibitem{sort2016}
Alex Bewley, Zongyuan Ge, Lionel Ott, Fabio Ramos, and Ben Upcroft.
\newblock Simple Online and Realtime Tracking.
\newblock In \emph{Proceedings of the IEEE International Conference on Image Processing (ICIP)}, 2016.
\newblock \url{https://arxiv.org/abs/1602.00763}.

\bibitem{tapir2023}
Carl Doersch, Yi Yang, Mel Vecerik, Dilara Gokay, Ankush Gupta, Yusuf Aytar, Joao Carreira, and Andrew Zisserman.
\newblock TAPIR: Tracking Any Point with Per-Frame Initialization and Temporal Refinement.
\newblock In \emph{Proceedings of the IEEE/CVF International Conference on Computer Vision (ICCV)}, 2023.
\newblock \url{https://arxiv.org/abs/2306.08637}.

\bibitem{deva2023}
Ho Kei Cheng, Seoung Wug Oh, Brian Price, Alexander Schwing, and Jia-Bin Huang.
\newblock Tracking Anything with Decoupled Video Segmentation.
\newblock In \emph{Proceedings of the IEEE/CVF International Conference on Computer Vision (ICCV)}, 2023.
\newblock \url{https://arxiv.org/abs/2309.03903}.

\bibitem{motiongs2024}
Ruijie Zhu, Zhihao Liang, Xingyu Chen, Jingyi Yu, and others.
\newblock Exploring Explicit Motion Guidance for Deformable 3D Gaussian Splatting.
\newblock In \emph{Advances in Neural Information Processing Systems (NeurIPS)}, 2024.
\newblock \url{https://proceedings.neurips.cc/paper_files/paper/2024/hash/b88318174aad2cc174a4e05ab6bfad80-Abstract-Conference.html}.

\bibitem{sam2_2024}
Nikhila Ravi, Valentin Gabeur, Yuan-Ting Hu, Ronghang Hu, Chaitanya Ryali, Tengyu Ma, Haitham Khedr, Roman R\"adle, Chloe Rolland, Laura Gustafson, Eric Mintun, Junting Pan, Kalyan Vasudev Alwala, Nicolas Carion, Chao-Yuan Wu, Ross Girshick, Piotr Doll\'ar, and Christoph Feichtenhofer.
\newblock SAM 2: Segment Anything in Images and Videos.
\newblock arXiv preprint arXiv:2408.00714, 2024.
\newblock \url{https://arxiv.org/abs/2408.00714}.

\bibitem{samurai2024}
Cheng-Yen Yang, Cheng-Han Lee, Chun-Yi Chiang, Yi-Ting Chen, and Winston H. Hsu.
\newblock SAMURAI: Adapting Segment Anything Model for Zero-Shot Visual Tracking with Motion-Aware Memory.
\newblock arXiv preprint arXiv:2411.11922, 2024.
\newblock \url{https://arxiv.org/abs/2411.11922}.

\bibitem{samite2025}
Qingze Xu, Fermin Segura, Yue Huang, and Luc Van Gool.
\newblock SAMITE: Position Prompted SAM2 with Calibrated Memory for Zero-Shot Visual Tracking.
\newblock arXiv preprint arXiv:2507.21732, 2025.
\newblock \url{https://arxiv.org/abs/2507.21732}.

\end{thebibliography}

\end{document}
